Sia $n\in\mathbb{N}$, definiamo $\mathbb{R}^n$ come \[\mathbb{R} := \underbrace{\mathbb{R}\times
\mathbb{R}\times \dots \times\mathbb{R}}_{n \text{ volte}} = \left\{\overline{x}:=\left(x_1,
x_2, \ldots , x_n\right)\mid \forall j\in\mathbb{N},\;x_j\in\mathbb{R}\right\}.\]

Definiamo inoltre le seguenti operazioni in $\mathbb{R}^n$.

\begin{definition}
    \textbf{Somma in $\mathbb{R}^n$}
    
    \[
        \overline{x} + \overline{y} := \left(x_1+y_1, x_2+y_2, \ldots, x_n+y_n\right).
    \]
\end{definition}
\begin{definition}
    \textbf{Prodotto per uno scalare}
    
    \[
        \forall \alpha\in\mathbb{R},\; \alpha\cdot\overline{x} :=\left(\alpha x_1,
        \alpha x_2, \ldots, \alpha x_n\right).
    \]
\end{definition}

Si può facilmente dimostrare che le due operazioni di somma e prodotto per scalare così
definite in $\mathbb{R}^n$ godono degli assiomi/proprietà caratteristiche di uno spazio
vettoriale (elencate in \defref{def:spazio-vettoriale}).
Possiamo dunque dire che $\mathbb{R}^n$ con le operaizoni deinfinite sopra è uno spazio vettoriale.

Nello spazio vettoriale così identificato possiamo introdurre un operazione, definita
come segue.

\begin{definition}
    \textbf{Prodotto scalare euclideo}
    
    \[
        \langle \overline{x}, \overline{y}\rangle := x_1 y_1 + x_2 y_2 + \ldots + x_n+y_n
         = \sum_{i = 1}^{n} x_i y_i.
    \]
\end{definition}

Si può facilmente mostrare che tale operazione rispetta tutte gli assiomi / proprietà
caratteristiche necessarie per un prodotto vettoriale.
Tale prodotto vettoriale prende il nome di prodotto vettoriale euclideo.

Dato dunque un prodotto vettoriale in uno spazio vettoriale, viene automaticamente indotta
la norma seguente, chiamata norma euclidea poiché indotta dal prodotto euclideo.



\begin{definition}
    \textbf{Norma euclidea}
    \[
        \| \overline{x} \| := \sqrt{\langle \overline{x}, \overline{x} \rangle}.
    \]
\end{definition}

